\section{The common marginal after tracing the good middle region}
\label{sec:replica_good_configuration_pair_marginal}

\begin{definition}[The five-factor regrouping]
\label{def:five_factor_pair_middle_equiv}
\lean{TensorPower.fiveFactorPairMiddleEquiv}
\leanok
Regroup the good-copy factors, with their fixed common enumeration, by
\begin{align*}
 (Q_G,Y_G,V_G,C_G,R_G)
 \longleftrightarrow ((Q_G,C_G),(V_G,R_G)),Y_G.
\end{align*}
Write $e$ for this coordinate equivalence. The ordering is the one used
for the paired merge estimate in
\cite[Section~8, lines~513--549]{OpenAI2026AreaLaw}.
\end{definition}

\begin{theorem}[The actual density and its common marginal]
\label{thm:replica_good_configuration_pair_marginal}
\lean{Matrix.posSemidef_replicaGoodConfigurationMarginal,Matrix.partialTraceRight_replicaGoodConfigurationMarginal_pairMiddle,Matrix.trace_replicaGoodConfigurationMarginal}
\leanok
\uses{def:replica_good_configuration_marginal,def:replica_good_pair_marginal}
Let $w$ be the actual excitation component for a subset $B$ of the
copies. Let $\rho$ be its reduced density retaining all five good
factors $Q_G,Y_G,V_G,C_G,R_G$, and tracing only the bad physical and
auxiliary coordinates. In the coordinates of $e$, denote this matrix
by $\rho_e$. Then $\rho$ is positive semidefinite,
\begin{align*}
 \operatorname{Tr}_{Y_G}\rho_e&=\rho_{\mathrm{pair}},\\
 \operatorname{Tr}\rho&=\|w\|^2,
\end{align*}
where $\rho_{\mathrm{pair}}$ is the actual common good-pair marginal.
Neither the original vector nor its component is required to be
normalized or permutation invariant. This is the common density used
in \cite[Section~8, lines~513--549]{OpenAI2026AreaLaw}.
\end{theorem}
\begin{proof}
\leanok
\uses{def:five_factor_pair_middle_equiv,thm:replica_good_pair_marginal_trace}
Positivity follows from positivity of the vector density and partial
trace. Both formulas for the common marginal sum the same matrix
entries over, in order, the good middle coordinates, the bad physical
coordinates, and the two bad auxiliary registers. The trace is preserved
by the coordinate change and partial trace, and equals the squared norm
of the original component.
\end{proof}

\begin{theorem}[Expectations of paired observables]
\label{thm:replica_good_configuration_pair_expectation}
\lean{Matrix.trace_replicaGoodConfigurationMarginal_pairMiddle_mul}
\leanok
For every matrix $H$ on the two good pairs,
\begin{align*}
 \operatorname{Tr}\bigl(\rho_e(H\otimes I_{Y_G})\bigr)
 =\operatorname{Tr}(\rho_{\mathrm{pair}}H).
\end{align*}
No positivity or Hermiticity of $H$ is required.
\end{theorem}
\begin{proof}
\leanok
\uses{thm:replica_good_configuration_pair_marginal}
Use the adjoint identity for partial trace and the preceding theorem.
\end{proof}
